% generator_zoo.tex -- the catalogue appendix.
%
% The longtable proper is auto-generated by
% src/utilities/generator_zoo.py from data/generator_catalogue.json
% and \input{}'d at the end of this section.  This file provides
% the hand-written reader's map, per-class structural notes, and
% pointers to the JSON schema.

\label{sec:generator_zoo}

The 71 generators of the per-site automorphism centralizer are
catalogued in Table~\ref{tab:generator_zoo}.  This appendix
describes the row schema, the seven SM-classification blocks the
table is partitioned into, and the JSON emit that machine-readable
downstream consumers should use rather than scraping the
\LaTeX{} table.

\subsection{Row schema}
\label{subsec:zoo_row_schema}

Each of the 71 rows of Table~\ref{tab:generator_zoo} records:

\begin{description}
\item[Name.]  A stable identifier of the form \texttt{G\_NN} with
  \texttt{NN} $\in \{01, \ldots, 71\}$.  Names are preserved across
  catalogue revisions and are the canonical reference for downstream
  papers.  The ordering convention is documented in the docstring
  of \texttt{src/utilities/generator\_zoo.py}; in brief: SM block
  first (\texttt{G\_01} = $J_1$, \texttt{G\_02--G\_04} = $SU(2)_W$,
  \texttt{G\_05--G\_12} = $SU(3)_c$), then extras in the order
  iso-colour mixings, $\sigma_x$-iso, $\sigma_x$-colour,
  $\sigma_x$-iso-col.
\item[Chir / Iso / Col.]  The tensor-signature factor of the
  generator on $\mathbb{C}^{12} = \mathbb{C}^2_{\text{chir}}
  \otimes \mathbb{C}^2_{\text{iso}} \otimes \mathbb{C}^3_{\text{col}}$.
  Possible values: $I$ (identity), $\sigma_x / \sigma_y / \sigma_z$
  (Paulis, for the chirality and isospin factors), or
  $\lambda_1, \ldots, \lambda_8$ (Gell-Manns, for the colour
  factor).  Only $I$ and $\sigma_x$ appear in the chirality column
  because the centralizer condition forces $[X, \sigma_x \otimes
  I_2 \otimes I_3] = 0$.
\item[Tag.]  A 3-bit factor-action tag $(c, i, k) \in \{0, 1\}^3$,
  one bit per factor, with bit value $1$ iff the corresponding
  tensor factor is non-identity.  Quick at-a-glance reading of
  which tensor factors a generator touches.
\item[$\mathbf{SU(3)}$.]  The colour irrep, either
  $\mathbf{8}$ (adjoint of $SU(3)_c$) or $\mathbf{1}$ (trivial).
\item[$\mathbf{SU(2)}$.]  The weak-isospin irrep, either
  $\mathbf{3}$ (adjoint of $SU(2)_W$) or $\mathbf{1}$ (trivial).
\item[Block.]  The chirality eigenspace ($I$ or $\sigma_x$) that
  the generator lives in.  In the structural decomposition
  $\mathfrak{Aut}_{\text{cent}} = \mathfrak{su}(6)_+ \oplus
  \mathfrak{su}(6)_- \oplus \mathfrak{u}(1)$, the $I$-block
  generators sit in $\mathfrak{su}(6)_+$ and the $\sigma_x$-block
  generators in $\mathfrak{su}(6)_-$ (with $J_1$ as the central
  $\mathfrak{u}(1)$).
\item[Class.]  The SM-classification subgroup the generator
  belongs to.  Seven possible values: \texttt{J\_1},
  \texttt{SU(2)\_W}, \texttt{SU(3)\_c} (the 12 SM generators);
  \texttt{iso\_col\_mixing} (24, leptoquark-flavoured),
  \texttt{sigmaX\_iso} (3, chirality-shadow $SU(2)_W$),
  \texttt{sigmaX\_col} (8, chirality-shadow $SU(3)_c$),
  \texttt{sigmaX\_iso\_col} (24, chirality-shadow leptoquark).
\end{description}

\subsection{The seven blocks at a glance}
\label{subsec:zoo_blocks}

\begin{description}
\item[\texttt{J\_1} (1 generator).]  The central $\mathfrak{u}(1)$
  of the centralizer.  $J_1 = \sigma_x \otimes I_2 \otimes I_3$
  commutes with everything in the 71-dim algebra and is the
  hypercharge-like central charge of the SM block.

\item[\texttt{SU(2)\_W} (3 generators).]  The standard weak-isospin
  generators $I_2 \otimes \sigma_a \otimes I_3$, $a \in \{x, y, z\}$.
  These are SM gauge.

\item[\texttt{SU(3)\_c} (8 generators).]  The standard colour
  Gell-Manns $I_2 \otimes I_2 \otimes \lambda_a$, $a \in \{1,
  \ldots, 8\}$.  These are SM gauge.

\item[\texttt{iso\_col\_mixing} (24 generators).]  Leptoquark-
  flavoured operators $I_2 \otimes \sigma_a \otimes \lambda_b$.
  Transform as $(\mathbf{8}, \mathbf{3})$ under the SM adjoint
  action -- colour-octet, weak-triplet, exactly the GUT
  leptoquark irrep.  In the chirality-$I$ block.

\item[\texttt{sigmaX\_iso} (3 generators).]  Chirality-shadow
  $SU(2)_W$ operators $\sigma_x \otimes \sigma_a \otimes I_3$.
  Same $(\mathbf{1}, \mathbf{3})$ SM-irrep content as
  $\mathfrak{su}(2)_W$, but living in the $\sigma_x$ chirality
  block, so they are not SM gauge.

\item[\texttt{sigmaX\_col} (8 generators).]  Chirality-shadow
  $SU(3)_c$ operators $\sigma_x \otimes I_2 \otimes \lambda_a$.
  Same $(\mathbf{8}, \mathbf{1})$ content as $\mathfrak{su}(3)_c$,
  in the $\sigma_x$ block.

\item[\texttt{sigmaX\_iso\_col} (24 generators).]  Chirality-shadow
  leptoquark operators $\sigma_x \otimes \sigma_a \otimes \lambda_b$.
  Same $(\mathbf{8}, \mathbf{3})$ content as the
  \texttt{iso\_col\_mixing} block, in the $\sigma_x$ block.
\end{description}

Sums: SM block $= 1 + 3 + 8 = 12$; extras block $= 24 + 3 + 8 + 24
= 59$; total $= 71$.  Of the 59 extras, $48 = 24 + 24$ are
leptoquark-flavoured (in $(\mathbf{8}, \mathbf{3})$) and the
remaining $11 = 3 + 8$ are chirality-shadow SM-flavoured (in
$(\mathbf{1}, \mathbf{3}) \oplus (\mathbf{8}, \mathbf{1})$).

\subsection{Machine-readable emit}
\label{subsec:zoo_json}

The catalogue is also emitted as
\texttt{data/generator\_catalogue.json}, with one record per
generator (same fields as Table~\ref{tab:generator_zoo}) plus a
full bracket-class table: for each of the
$\binom{71}{2} = 2485$ unordered pairs $(G_i, G_j)$ with $i < j$,
the class of the bracket $[G_i, G_j]$ (one of \texttt{zero} /
\texttt{in\_SM} / \texttt{in\_extras} / \texttt{mixed}) is
recorded.  Downstream tools that consume the catalogue
programmatically should read the JSON and refer to generators
by their \texttt{G\_NN} name; this LaTeX longtable is the human-
facing companion.

The JSON schema is documented in the docstring of
\texttt{src/utilities/generator\_zoo.py}.  Regenerate the JSON
and this longtable in lockstep via:

\begin{verbatim}
python -m src.utilities.generator_zoo
\end{verbatim}

\subsection{Bracket-class summary}
\label{subsec:zoo_brackets}

The full bracket-class table is in the JSON emit; here we
record the aggregate counts (also reproduced by
\texttt{src/utilities/aut\_centralizer\_extras\_commutators.py}
for the subsets it covers):

\begin{itemize}
\item $[\text{SM}, \text{SM}]$ (66 pairs): closure of the 12-dim
  SM gauge subalgebra inside the centralizer.  All brackets land
  in SM (the SM block is a Lie subalgebra).
\item $[\text{extras}, \text{SM}]$ (708 pairs): 256 zero, 0
  in\_SM, 452 in\_extras, 0 mixed.  The 59 extras form an
  invariant module under the adjoint action of SM.
\item $[\text{extras}, \text{extras}]$ (1711 pairs): 339 zero,
  178 in\_SM, 1146 in\_extras, 48 mixed.  The extras subspace is
  \emph{not} a Lie ideal; brackets within it produce SM
  components, consistent with the standard algebraic shape of a
  broken-symmetry pattern (see
  \texttt{notes/aut\_centralizer\_extras\_commutators.md}).
\end{itemize}

Sum: $66 + 708 + 1711 = 2485 = \binom{71}{2}$.

\subsection{The catalogue}
\label{subsec:zoo_catalogue}

%% The longtable below is auto-generated by
%% src/utilities/generator_zoo.py.  Do not edit by hand; regenerate
%% via 'python -m src.utilities.generator_zoo'.  Wrapped in
%% \IfFileExists{} so a fresh clone whose user has not yet run
%% the catalogue script still builds; the placeholder branch then
%% prints a visible reminder.
\IfFileExists{sections/generator_zoo_table.tex}{%
    \input{sections/generator_zoo_table}%
}{%
    \placeholder{Auto-generated catalogue table not present.
    Regenerate via \texttt{python -m src.utilities.generator\_zoo}
    from the repo root; the script writes
    \texttt{paper/sections/generator\_zoo\_table.tex} and
    \texttt{data/generator\_catalogue.json} together.}%
}
